\section{Introduction}
\label{sec:introduction}
% --- P1: VLAs are strong but spatially fragile ---
Vision-Language-Action (VLA) models~\cite{rt22023arxiv, octo_2023, kim2024openvla, pi0, pi05, gr00tn1_2025} inherit rich semantic representations from large-scale vision-language pretraining, enabling them to follow free-form instructions and recognize a wide range of objects during manipulation. However, VLAs remain brittle in test-time generalization along two distinct axes—\emph{spatial generalization}, where object positions differ from those seen during training~\cite{zhou2025liberopro, fei25libero-plus}, and instance-level \emph{task generalization}, where a familiar scene is paired with an instruction specifying a different target at test time~\cite{zhou2025liberopro}. Closing these generalization gaps has become a central challenge for the field.

% --- P2: Current solution: spatial grounding ---
A growing line of work addresses this by providing the policy with explicit spatial grounding---the location of the task-relevant object or placement region. One family relies on external grounding modules such as Vision-Language Models (VLMs) or detectors~\cite{yuan2024robopoint, roboground} and conveys their predictions to the policy as coordinates in the language prompt~\cite{molmobot}, visual prompt~\cite{moka, wang2026vp}, or combined visual--textual conditioning~\cite{coavla}. The other trains the VLA itself to emit such spatial outputs as an intermediate chain-of-thought prediction before producing actions~\cite{molmoact2025, ecot, steerablepolicy}. While both directions improve generalization, neither systematically addresses how the spatial signal should be represented and injected into the policy to achieve robust generalization.

% --- P3: Our key insight ---
% Specifically, we show that, by lifting a 2D point through depth into 3D and injecting the 3D point directly into the action head is a more effective way to unlock robust generalization, rather than relaying it through 2D proxies such as language tokens or image overlays. We validate this through controlled experiments: when all methods receive an identical oracle target point, encoding it as textual coordinates in the language prompt or rendering it as a visual prompt yields only modest improvements, whereas lifting it into 3D and injecting it as a spatial embedding produces dramatic gains. 
\begin{wrapfigure}{r}{0.45\textwidth}
    \includegraphics[width=\linewidth]{assets/figures/teaser_fig.pdf}
    \caption{\textbf{Comparison of grounding representation and injection paradigms.} Given an identical 2D target point, appending coordinates to the language prompt and overlaying it as a visual prompt on the image yield slight improvements over the baseline, whereas lifting it into 3D and injecting it into the action head leads to substantially better performance on LIBERO-PRO benchmark.} 
    \label{fig:teaser}
\end{wrapfigure}

Specifically, as shown in Fig.~\ref{fig:teaser}, we show that lifting a 2D point through depth into 3D and injecting it directly into the action head as a spatial embedding is more effective at unlocking robust generalization than relaying it through 2D proxies such as language tokens or image overlays. 
The reason is intuitive: action prediction operates in the 3D space, so a 3D signal matches the geometry of the output, whereas 2D representations force the policy to internally reconstruct this geometry from visual and linguistic features. The same mechanism also explains improvements on task generalization: an explicit spatial signal tied to the current instruction prevents the policy from defaulting to scene-triggered trajectories memorized during training, forcing it to act on the target specified by the instruction at test time.

%Specifically, we show that, given an adequate 2D grounding signal, lifting the signal into 3D and injecting it directly into the action head is a more effective way to unlock robust generalization, rather than relaying it through 2D proxies such as language tokens or image overlays. We validate this through controlled experiments: when all methods receive an identical oracle target point, encoding it as textual coordinates in the language prompt or rendering it as a visual marker yields only modest improvements, whereas lifting it into 3D and injecting it as a spatial embedding produces dramatic gains. The reason is intuitive: action prediction operates in 3D physical space, so a 3D signal matches the geometry of the output, whereas 2D representations force the policy to internally reconstruct this geometry from visual and linguistic features. The same mechanism also explains improvements on task generalization: an explicit spatial signal tied to the current instruction prevents the policy from defaulting to scene-triggered trajectories memorized during training, forcing it to act on the target specified by the instruction at test time.

% --- P4: Our method ---
To realize this design, we propose a lightweight, model-agnostic module that lifts the grounding signal into 3D and feeds it into the action head. Specifically, we lift the 2D target point into 3D using depth and camera parameters and compute its displacement relative to the gripper. We then use a two-layer MLP to encode this displacement into a spatial embedding. Modern VLAs commonly adopt a diffusion transformer (DiT) as the action head~\cite{li2024cogact, pi0, pi05, gr00tn1_2025, molmobot}; we add the spatial embedding to the existing time-step embedding, and the combined embedding conditions the DiT via adaptive layer normalization~\cite{DiT, perez2018film}. The module requires no changes to the VLA backbone or pretraining pipeline; the only added component is a two-layer MLP.

% --- P5: Results ---
% TODO: add results paragraph when experimental data is more complete
We validate this design through extensive simulation experiments on LIBERO~\cite{libero} and LIBERO-PRO~\cite{zhou2025liberopro}, complemented by real-world evaluations. Applied on top of GR00T-N1.6, our module raises the average LIBERO-PRO success rate from $31.2\%$ to \textbf{77.5}$\%$ under Task Perturbation and from $28.1\%$ to \textbf{60.2}$\%$ under Position Perturbation. Our method yields comparable gains when applied to $\pi_{0.5}$ ($37.3\%$ to \textbf{75.9}$\%$ and $47.0\%$ to \textbf{72.2}$\%$ under the two perturbations), demonstrating that the mechanism is backbone-agnostic across VLAs with diffusion-based action heads. 
We further validate our method in the real world, where it remains effective under both perturbations and substantially outperforms the baselines.
These results point to a simple but overlooked finding: with adequate 2D grounding in hand, generalization hinges on how the signal is represented and injected—lifting it into 3D and injecting it directly into the action head substantially improves both spatial and instance-level task generalization.